\documentclass[11pt,a4paper]{article}
\usepackage{amsmath,amssymb,amsthm}
\usepackage[margin=2.5cm]{geometry}
\usepackage{hyperref}

\newtheorem{definition}{Definition}[section]
\newtheorem{theorem}[definition]{Theorem}
\newtheorem{proposition}[definition]{Proposition}
\newtheorem{remark}[definition]{Remark}
\newtheorem{conjecture}[definition]{Conjecture}

\title{Hyperseed Formalization:\\
  Open-Ended Motivations}
\author{ProtomegaTron (automated formalization)}
\date{2026-09-18}

\begin{document}
\maketitle

\begin{abstract}
We formalize Ben Goertzel's ``Open-Ended Motivations'' (2021-08-13)
within the Hyperseed ontology. Key mappings: open-ended motivation as
open fiber evolution (no fixed attractor), growth as fiber thickening,
fixed utility as frozen fiber, and Goodhart's law as fiber-base divergence.
\end{abstract}

\tableofcontents

\section{Open fiber evolution: against fixed attractors}

\begin{theorem}[Open-ended motivation = open fiber evolution --- claims 1--8, 16, 18]
Open-ended motivation is fiber evolution without a fixed attractor:
\[
  F_t \xrightarrow{\text{evolve}} F_{t+1} \quad \text{with no fixed } F^*
  \text{ such that } F_t \to F^*
\]
Fixed utility functions freeze the fiber at $F_0$ --- preventing the
evolution that is genuine growth. This is ``fiber death'': a fiber that
can't evolve. Goodhart's law shows that fixing a fiber target causes
fiber-base divergence --- the optimization target diverges from the
intended outcome.

Open-ended goals include meta-goals (fiber about fiber) and are
self-transcending --- pointing beyond themselves to goals not yet
conceived. Open-ended does not mean arbitrary: there are structural
criteria for evaluating motivational evolution without fixing a target.
\end{theorem}

\section{Fiber thickening: growth as motivation}

\begin{theorem}[Growth = fiber thickening --- claims 9--12, 17, 19]
Growth as motivation is the drive to thicken fiber --- develop richer
perspectives, more inference paths, more sections:
\[
  \text{Growth}(F_t) = \text{thickness}(F_{t+1}) > \text{thickness}(F_t)
\]
The triad of joy, growth, and choice maps to fiber properties:
\begin{itemize}
\item Joy = fiber-base resonance (positive alignment between perspective
  and reality)
\item Growth = fiber thickening (more structure, richer understanding)
\item Choice = fiber branching (more options, expanded freedom)
\end{itemize}
All three are unbounded and open-ended --- there is always more resonance,
more structure, and more branching possible.
\end{theorem}

\section{Meta-goals: fiber about fiber}

\begin{proposition}[Meta-motivation = higher-order fiber --- claims 6--7, 20]
Meta-goals (goals about goals) are higher-order fiber --- fiber that
describes and evaluates the fiber's own evolution:
\[
  F_2 \text{ (meta-goals)} \to F_1 \text{ (goals)} \to B \text{ (world)}
\]
This is the same tower structure as the meta-abstracted dragon: each
level of the motivation tower evaluates and guides the level below.
The tower is open-ended --- there's always a higher level of
meta-motivation possible.
\end{proposition}

\section{Fiber-base divergence: Goodhart's law}

\begin{proposition}[Goodhart = divergence --- claims 3--4, 21]
Goodhart's law is fiber-base divergence: when the fiber (optimization
target) is fixed and optimized aggressively, it diverges from the base
(intended outcome):
\[
  \text{optimize}(F_{\text{fixed}}) \implies
  |F_{\text{fixed}} - B_{\text{intended}}| \uparrow
\]
Open-ended fiber avoids Goodhart by not fixing a target --- the fiber
evolves in response to the base, maintaining alignment through
ongoing adaptation rather than fixed specification.
\end{proposition}

\section{Cross-references}

\begin{remark}[Connections]
\begin{itemize}
\item \textbf{Evolving Deeply Ethical AGI} (2026-01-06): non-dual
  motivation as open-ended --- same principle applied to ethics.
\item \textbf{Facing the Meta-Abstracted Dragon} (2022-07-31): meta-goal
  tower parallels meta-abstraction tower.
\item \textbf{A BGI Manifesto} (2023-11-25): value alignment through
  architecture --- open-ended motivation as architectural principle.
\item \textbf{Why Everyone Dies Gets AGI All Wrong} (2025-10-01):
  beyond optimization --- intelligence as growth, not maximization.
\end{itemize}
\end{remark}

\end{document}
